\section{Related Work}
\label{sec:related_work}

\textbf{Parameter-Efficient Fine-Tuning (PEFT).} PEFT has emerged as the dominant paradigm for adapting massive pre-trained foundation models to downstream applications. Among these, Low-Rank Adaptation (LoRA) \cite{hu2021lora} is particularly prominent. LoRA factorizes weight updates into two low-rank matrices, $A \in \mathbb{R}^{D \times r}$ and $B \in \mathbb{R}^{r \times D}$, significantly reducing the memory footprint during training while preserving representational capacity. Other PEFT methods, such as adapters \cite{houlsby2019parameter}, prefix tuning \cite{li2021prefix}, and prompt tuning \cite{lester2021power}, follow similar modular architectures. Although highly effective, serving $K$ parallel experts independently leads to linear memory scaling and high serving costs, prompting researchers to seek unified multi-task ensembling alternatives.

\textbf{Static Model Merging.} Statically combining model weights is a popular strategy for consolidating multiple task experts into a single set of parameters. Techniques such as Task Arithmetic \cite{ilharco2022editing}, TIES-Merging \cite{yadav2023ties}, and DARE \cite{dares2023} resolve parameter-level interference (e.g., redundant parameter signs and magnitudes) and average weight vectors across experts. While statically merged models have zero extra inference overhead, they suffer from two fundamental limitations: (1) capacity dilution, where merging conflicting task vectors washes out individual task performance, and (2) static rigidity, as they are unable to adapt to mixed, heterogeneous query streams at test time.

\textbf{Mixture of Experts (MoE) and Routing.} Sparse Mixture of Experts (MoE) \cite{shazeer2017outrageously, fedus2022switch} scales network capacity by activating a sparse subset of expert pathways for each token or sample. However, standard MoE models rely on parameteric gating networks that must be trained jointly with the model, which is prone to routing collapse and requires heavy optimization. Moreover, they are difficult to adapt to post-hoc ensembling of independently fine-tuned PEFT experts.

\textbf{Dynamic Model Merging and Serving.} To address the limitations of static merging and parametric MoEs, recent works have proposed dynamic serving and merging techniques. Frameworks like AdaMerging \cite{adamerging2023} and SPS \cite{sps2024} dynamically compute merging coefficients based on input characteristics. However, they rely on trainable routing modules. Recent state-of-the-art frameworks have attempted to make routing training-free, but they have introduced substantial systems and algorithmic complexity:
\begin{itemize}
    \item \textbf{SPS-ZCA} \cite{sps_zca2025} performs zero-shot routing but relies on high-dimensional centroids pre-computed from a 64-sample calibration split, Unit-Norm Calibration (UNC), Intra-Task Dispersion Calibration (IDC) scaling factors, and multi-dimensional Expectation-Maximization-fitted Gaussian Mixture Models (GMMs) for out-of-distribution (OOD) rejection.
    \item \textbf{SABLE} \cite{sable2025} and \textbf{PFSR} \cite{pfsr2025} route inputs by computing cosine similarities between early activations and frozen classification-head weights.
\end{itemize}

These head-dependent routers have major limitations. First, they cannot run in settings where classification heads do not exist (e.g., autoregressive decoder layers, or head-free embedding extraction). Second, they suffer from the \emph{Early-Layer Routing Paradox}, where task-specific adapters must run in parallel throughout the entire depth of the model before classification-head similarity can be computed. This causes severe representational interference and activation dilution in early layers.

\textbf{Multi-Tenant Serving Systems.} In the systems literature, executing multiple heterogeneous LoRA experts simultaneously has been heavily optimized using frameworks like S-LoRA \cite{slora2024} and Punica \cite{punica2024}. These frameworks design specialized GPU kernels (such as \texttt{bgmv}) to execute batch-level, multi-tenant low-rank matrix-vector multiplications, completely bypassing sequential parameter reloads and weight merging. While S-LoRA and Punica represent the state-of-the-art for GPU servers, they focus entirely on hardware-level execution parallelization and do not address the mathematical problems of dynamic routing, ensembling, or OOD rejection. LSPR is highly complementary to these systems: it provides a clean, zero-overhead mathematical routing and activation-space blending layer, making it an ideal choice for resource-constrained edge CPUs where specialized multi-tenant CUDA kernels are completely unsupported and sequential DRAM loading is the default native alternative.

In contrast to these complex, over-engineered pipelines, our proposed LoRA Subspace Projection Routing (LSPR) is training-free and requires zero task-specific calibration data at deployment time, while being completely head-free and hyperparameter-light. LSPR leverages the closed-form QR decomposition of the low-rank PEFT matrices themselves to achieve elegant, scale-invariant, sample-specific dynamic routing at early layers, establishing a new SOTA on both accuracy and systems latency.
